% This document was generated by an LLM.

\documentclass{essay}

\newtheorem*{theorem*}{Theorem}

\title{The Banach Fixed-Point Theorem:\\
From a Calculator Trick to Spaces of Functions}
\author{}
\date{}

\begin{document}

\maketitle

\section{A calculator puzzle}

Take any calculator, type in a number, and repeatedly press the cosine
button, working in radians. Start, say, with $x_0 = 1$, and compute
$x_1 = \cos(x_0)$, then $x_2 = \cos(x_1)$, and so on. No matter what
number you start with, the sequence converges to the same value,
\[
  x \approx 0.739085\ldots,
\]
which is the unique solution of the equation
\[
  x = \cos(x).
\]
Such a number is called a \emph{fixed point} of the function
$f(x) = \cos(x)$: feeding it into $f$ returns the same number. The
puzzle is why this always happens, regardless of the starting point,
and what property of the cosine function is responsible. The answer
is the Banach fixed-point theorem, and building it up from ordinary
single-variable calculus is the purpose of this essay. Along the way
the same argument will be shown to generalize far beyond the real
line, to the point where it underwrites the existence and uniqueness
of solutions to differential equations.

\section{Contractions and the Mean Value Theorem}

Suppose $f$ is differentiable and its derivative satisfies
$|f'(x)| \le k$ for some constant $k < 1$ on some interval. The Mean
Value Theorem says that for any two points $x, y$ in that interval
there is a point $c$ between them with
\[
  f(x) - f(y) = f'(c)(x - y),
\]
and hence
\[
  |f(x) - f(y)| \le k\,|x - y|.
\]
In words: $f$ shrinks the distance between any two points by a factor
of at least $k$. A function with this property is called a
\emph{contraction}.

The cosine function is a contraction near its fixed point. Its
derivative is $f'(x) = -\sin(x)$, so $|f'(x)| \le 1$ everywhere, and
in fact $|f'(x)| < 1$ away from the origin. Near the fixed point
$0.739\ldots$, one has $\sin(x) \approx 0.67$, comfortably below $1$.
This is the entire mechanism behind the calculator experiment:
cosine contracts distances, and contractions force convergence, as
the next section shows.

\section{Why a contraction forces convergence}

Define the iteration $x_{n+1} = f(x_n)$ starting from an arbitrary
point $x_0$, where $f$ is a contraction with constant $k < 1$. Look
at the distance between consecutive terms:
\[
  |x_2 - x_1| = |f(x_1) - f(x_0)| \le k\,|x_1 - x_0|,
\]
\[
  |x_3 - x_2| = |f(x_2) - f(x_1)| \le k\,|x_2 - x_1| \le k^2\,|x_1 - x_0|,
\]
and in general
\[
  |x_{n+1} - x_n| \le k^n\,|x_1 - x_0|.
\]
The steps between successive terms shrink geometrically. This is the
familiar situation from geometric series: if the step sizes decay
like $k^n$ with $k < 1$, then the total distance still to be
travelled from step $n$ onward is bounded by a convergent geometric
series,
\[
  |x_{n+m} - x_n|
  \le \sum_{i=n}^{n+m-1} |x_{i+1} - x_i|
  \le k^n\,|x_1 - x_0| \sum_{i=0}^{\infty} k^i
  = \frac{k^n}{1-k}\,|x_1 - x_0|,
\]
and the right-hand side tends to $0$ as $n \to \infty$. The tail of
the sequence is therefore squeezed into an arbitrarily small
interval: the sequence has nowhere left to wander, so it must
converge to some limit $L$. Since $f$ is continuous, taking the limit
of both sides of $x_{n+1} = f(x_n)$ gives $L = f(L)$, so the limit is
itself a fixed point.

Uniqueness follows from the same contraction inequality. If $L_1$ and
$L_2$ were both fixed points, then
\[
  |L_1 - L_2| = |f(L_1) - f(L_2)| \le k\,|L_1 - L_2|,
\]
which, since $k < 1$, is only possible if $|L_1 - L_2| = 0$. So a
contraction has at most one fixed point, and the iteration finds it
starting from anywhere. The geometric-series bound above is not just
a convergence proof; it is also an explicit error estimate, giving
the number of iterations needed to approximate the fixed point to any
desired tolerance without knowing the fixed point in advance.

\section{Generalizing beyond the real line}

Nothing in the argument of the previous section used the fact that
the $x_n$ were real numbers. The only ingredients were:
\begin{itemize}
  \item a notion of distance between two points, written $|x - y|$
    on the real line, satisfying the usual triangle inequality;
  \item the fact that a sequence whose terms get arbitrarily close
    together must converge to a point that actually belongs to the
    space in question, a property known as \emph{completeness} (the
    real numbers have it; the rational numbers famously do not,
    since the decimal approximations to $\sqrt{2}$ get arbitrarily
    close together without converging to a rational number);
  \item a map $f$ that shrinks distances by a fixed factor $k < 1$.
\end{itemize}
None of this requires the points to be numbers. They could equally
well be points in $\mathbb{R}^n$, with distance measured by the
Euclidean norm and the contraction condition checked through the
Jacobian rather than a single derivative. They could even be entire
\emph{functions}. As long as the collection of objects in question
comes equipped with a sensible notion of distance, and is complete
with respect to that distance, the same telescoping-series argument
goes through unchanged. This is the content of the general theorem.

\begin{theorem*}[Banach fixed-point theorem]
Let $(X, d)$ be a complete metric space and let $f : X \to X$ satisfy
\[
  d(f(x), f(y)) \le k\, d(x, y)
\]
for some fixed constant $k < 1$ and all $x, y \in X$. Then $f$ has
exactly one fixed point in $X$, and for any starting point
$x_0 \in X$, the sequence defined by $x_{n+1} = f(x_n)$ converges to
that fixed point.
\end{theorem*}

The abstraction is not idle. It is exactly what is needed to make
sense of iteration on objects that are not numbers at all, and the
most important example of such an object is a function.

\section{The space of functions}

Consider two continuous functions $\varphi(x)$ and $\psi(x)$ defined
on an interval $[a, b]$. Their distance can be defined as the largest
vertical gap between their graphs,
\[
  d(\varphi, \psi) = \max_{x \in [a,b]} |\varphi(x) - \psi(x)|.
\]
This is a perfectly legitimate notion of distance: two functions are
close if their graphs stay close together everywhere on the
interval, and far apart if their graphs diverge substantially
anywhere, even at a single point. Equipped with this distance, the
collection of continuous functions on $[a,b]$ becomes a complete
metric space in exactly the sense required by the theorem: it has a
notion of distance, and a sequence of functions that get uniformly
closer and closer together does converge to another continuous
function on $[a,b]$. This last fact plays the same role that
completeness of the real numbers plays in ordinary calculus; it can
be taken on faith at this stage, in the same way that the
Intermediate Value Theorem is used without re-deriving the
completeness of $\mathbb{R}$ each time.

With this notion of distance in hand, a map that takes a function as
input and returns another function as output can be a contraction in
exactly the same sense as before, and the fixed point of such a map
is not a number but a function.

\section{Application: existence and uniqueness of solutions to ordinary differential equations}

This machinery is precisely what underlies the Picard--Lindel\"of
theorem. Consider an initial value problem
\[
  y' = f(x, y), \qquad y(x_0) = y_0.
\]
By the Fundamental Theorem of Calculus, this is equivalent to the
integral equation
\[
  y(x) = y_0 + \int_{x_0}^{x} f(t, y(t))\, dt.
\]
Define an operator $T$ that takes a candidate function $\varphi$ and
returns a new function $T\varphi$, given by
\[
  (T\varphi)(x) = y_0 + \int_{x_0}^{x} f(t, \varphi(t))\, dt.
\]
A solution of the original differential equation is precisely a
function $y$ satisfying $Ty = y$: a fixed point of $T$, in the sense
of the previous section, except that the ``point'' being sought is
an entire curve rather than a single number.

Suppose $f(x, y)$ satisfies a Lipschitz condition in its second
argument, meaning there is a constant $L$ such that
$|f(x, \varphi) - f(x, \psi)| \le L\,|\varphi - \psi|$ for all
relevant $\varphi, \psi$. Then for two candidate functions $\varphi$
and $\psi$,
\[
  |(T\varphi)(x) - (T\psi)(x)|
  = \left| \int_{x_0}^{x} \big[ f(t, \varphi(t)) - f(t, \psi(t)) \big]\, dt \right|
  \le \int_{x_0}^{x} L\, |\varphi(t) - \psi(t)|\, dt.
\]
Since $|\varphi(t) - \psi(t)| \le d(\varphi, \psi)$ for every $t$ in
the interval, by definition of the sup-norm distance, the integral on
the right is bounded by $L\,|x - x_0|\, d(\varphi, \psi)$. Restricting
attention to a sufficiently small interval of width $h$ around $x_0$,
chosen so that $Lh < 1$, gives
\[
  d(T\varphi, T\psi) \le \underbrace{Lh}_{\,k \,<\, 1\,} d(\varphi, \psi).
\]
This is exactly the contraction inequality from Section 2, transported
from the real line to the space of continuous functions. Everything
that followed from it there follows here as well. Iterating
$\varphi_{n+1} = T\varphi_n$ from any starting function, typically the
constant function $\varphi_0(x) = y_0$, produces a sequence of
functions whose distances shrink geometrically; by the argument of
Section 3, this sequence converges to a function $y$ satisfying
$Ty = y$, which is the unique solution of the initial value problem
on the chosen interval. The successive iterates $\varphi_0, \varphi_1,
\varphi_2, \ldots$ are precisely the Picard iterates, and the
geometric-series bound gives an explicit estimate of how closely a
given iterate approximates the true solution.

\section{Concluding remarks}

The path from a calculator experiment with cosine to the existence
and uniqueness of solutions of differential equations passes through
exactly one idea, applied twice. A contraction shrinks distances by a
fixed factor less than one; iterating a contraction produces a
sequence whose steps shrink geometrically; and a geometric sequence
of shrinking steps has nowhere to go but toward a single point, which
must be the unique fixed point of the map. On the real line this
point is a number, found by an idle sequence of button presses. On
the space of continuous functions, equipped with the natural notion
of distance between two graphs, the same argument produces a
function, and that function is the solution to a differential
equation. The abstraction from numbers to functions costs nothing in
the proof and gains a great deal in reach: the same page of reasoning
that explains why repeated cosine always lands on $0.739\ldots$ also
explains why a well-posed initial value problem has exactly one
solution.

\end{document}
